Reasoning lies at the core of intelligence, enabling logical inference, problem-solving, and decision-making across interactive and dynamic settings. Large language models (LLMs) have achieved remarkable gains in closed-world domains such as mathematical problem solving and code generation. Empirically, techniques that explicitize intermediate reasoning, such as Chain-of-Thought prompting, decomposition, and program-aided solving, have significantly bolstered inference performance \cite{wei2022chain,zhou2022least,gao2023pal,yao2023tree}. Yet, these approaches often assume static contexts and short-horizon reasoning. Conventional LLMs lack mechanisms to act, adapt, or improve in open-ended environments where information evolves over time.

In this survey, we systematize this evolution under the framework of \emph{Agentic Reasoning}: rather than passively generating sequences, LLMs are reframed as autonomous reasoning agents that plan, act, and learn through continual interaction with their environment. This reframing unifies \emph{reasoning} with \emph{acting}, positioning reasoning as the organizing principle for perception, planning, decision, and verification. Systems such as ReAct~\cite{yao2023react} interleave deliberation with environment interaction, tool-use frameworks enable self-directed API calling, and workflow-based agents dynamically orchestrate sub-tasks and verifiable actions~\cite{yao2023react,schick2023toolformer,shen2023hugginggpt}. Conceptually, this parallels the shift from static, one-shot inference to sequential decision-making under uncertainty. Unlike simple input-output mapping, this paradigm requires agents to plan over long horizons, navigate partial observability, and actively improve through feedback~\cite{wang2024survey,singh2025agentic,huang2024survey}.

\begin{tcolorbox}[
  agentdef,
  float,
  floatplacement=t
]
\textbf{Agentic reasoning} positions reasoning as the central mechanism of intelligent agents, spanning \emph{foundational capabilities} (planning, tool use, and search), \emph{self-evolving adaptation} (feedback, and memory-driven adaptation), and \emph{collective coordination} (multi-agent collaboration), realizable through either \emph{in-context} orchestration or \emph{post-training} optimization.
\end{tcolorbox}

To systematically characterize the environmental dynamics, we structure our survey around three complementary scopes of agentic reasoning: foundational capabilities, self-evolution, and collective intelligence, spanning diverse interactive and dynamic settings. \emph{\textbf{Foundational Agentic Reasoning}} establishes the bedrock of core single-agent capabilities, including planning, tool use, and search, that enable operations within stable, albeit complex, environments. Here, agents act by decomposing goals, invoking external tools, and verifying results through executable actions. For instance, program-aided reasoning~\cite{gao2023pal} grounds logical derivations in code execution; repository-level systems such as OpenHands~\cite{wang2024openhands} integrate reasoning, planning, and testing into unified loops; and structured memory modules~\cite{chhikara2025mem0,li2025memos} transform factual recall into procedural competence by persisting intermediate reasoning traces for reuse.

Building upon these foundations, \emph{\textbf{Self-Evolving Agentic Reasoning}} enables agents to improve continually through cumulative experience. Encompassing task-specific \emph{self-improvement} (e.g., via iterative critique), this paradigm extends adaptation to include persistent updates of internal states like memory and policy. Rather than following fixed reasoning paths, agents develop mechanisms for feedback integration and memory-driven adaptation to navigate evolving environments. Reflection-based frameworks such as Reflexion~\cite{shinn2023reflexion} allow agents to critique and refine their own reasoning processes, while reinforcement formulations such as RL-for-memory~\cite{yan2025memoryr1} formalize memory writing and retrieval as policy optimization. Through these mechanisms, agents dynamically integrate inference-time reasoning with learning, progressively updating internal representations and decision policies without full retraining. This continual adaptation links reasoning with learning, enabling models to accumulate competence, and generalize across tasks.

Finally, \emph{\textbf{Collective Multi-Agent Reasoning}} scales intelligence from isolated solvers to collaborative ecosystems. Rather than operating in isolation, multiple agents coordinate to achieve shared goals through explicit role assignment (e.g., manager–worker–critic), communication protocols, and shared memory systems~\cite{chen2023autoagents,hong2024metagpt}. As agents specialize in subtasks and refine each other’s outputs, collaboration amplifies reasoning diversity, enabling systems to debate, resolve disagreements, and achieve consistency through natural language-based multi-turn interactions~\cite{unleashing2024,wang2024battleagentbench}. However, this complexity also introduces challenges in stability, communication efficiency, and trustworthiness, necessitating structured coordination frameworks and rigorous evaluation standards~\cite{xiao2023agentbench,zhu2025multiagentbench}.


Across all layers, we analyze system constraints and optimization settings by distinguishing two complementary modes, corresponding to inference-time orchestration \cite{yao2023react,shinn2023reflexion,ni2024tree,li2025search,xu2025mem,wei2025evo} and training-based capability optimization \cite{ma2024coevolving,jin2025search,wei2025webagent,yan2025memoryr1}. \emph{\textbf{In-context Reasoning}} focuses on scaling inference-time compute: through structured orchestration, search-based planning, and adaptive workflow design, it enables agents to navigate complex problem spaces dynamically without modifying model parameters. Conversely, \emph{\textbf{Post-training Reasoning}} targets capability internalization: it consolidates successful reasoning patterns or tool-use strategies into the model's weights via reinforcement learning and fine-tuning. Together, they provide an actionable roadmap for designing agents.


\begin{tcolorbox}[
  agentscope,
  float,
  floatplacement=t
]
This survey reviews \emph{reasoning-empowered agentic systems} where reasoning drives adaptive behavior. We analyze these systems through two complementary optimization modes:
\begin{itemize}
    \item \textbf{In-context Reasoning}: scales inference-time interaction through structured orchestration and planning without parameter updates.
    \item \textbf{Post-training Reasoning}: internalizes reasoning strategies into model parameters via reinforcement learning and fine-tuning.
\end{itemize}
Our scope covers methodologies embedding these modes into planning, memory, and self-improvement across single-agent and multi-agent contexts. This survey summarizes progress up to 2025.
\end{tcolorbox}


Building on the three-layer taxonomy, agentic reasoning has begun to underpin a wide range of practical applications, from mathematical exploration \cite{trinh2024solving,romera2024mathematical} and vibe coding \cite{wang2024openhands,sapkota2025vibe,karpathy2025vibecoding} to scientific discovery \cite{bran2023chemcrowaugmentinglargelanguagemodels,bousetouane2025physicalaiagentsintegrating,ding2024matexpertdecomposingmaterialsdiscovery}, embodied robotics \cite{wang2023voyager,meghan2024embodiedrag,zhao2025embodied}, healthcare \cite{li2024mmedagent,huang2025biomni}, and autonomous web exploration~\cite{li2025websailor,zheng2025skillweaver}.
These applications expose distinct reasoning demands shaped by domain-specific data modalities, interaction constraints, and feedback loops, motivating diverse system designs \cite{sapkota2025ai,liu2024dynamic} that integrate planning, tool use, search, reflection, memory mechanisms, and multi-agent coordination.
On the other hand, the benchmark landscape has emerged to evaluate agentic reasoning, ranging from targeted tests that isolate individual agentic capabilities to application-specific benchmarks that assess end-to-end behavior in domain-specific environments and scenarios~\cite{shuyan2023webarena,jing2024visualwebarena,lawrence2024videowebarena,shridhar2020alfworld,xiao2023agentbench,zhu2025multiagentbench,deng2023mind2web,gou2025mind2web2}.


Together, this survey synthesizes agentic reasoning methods into a unified roadmap that bridges reasoning and acting. We systematically characterize these methods across the complementary scopes of foundational, self-evolving, and collective reasoning, while distinguishing between in-context and post-training optimization modes. We further contextualize this roadmap through representative applications and evaluation benchmarks, illustrating how different agentic reasoning mechanisms are instantiated and assessed across realistic domains and task settings. Finally, we outline open challenges and future directions, identifying key frontiers such as personalization, long-horizon interaction, world modeling, scalable multi-agent training, and governance frameworks for real-world deployment.





\begin{tcolorbox}[agentcontrib]
This survey makes the following contributions:
\begin{itemize}
\item \textbf{Conceptual framing}: We formalize the paradigm of \emph{Agentic Reasoning}, spanning foundational, self-evolving, and collective reasoning layers.
\item \textbf{Systematic review}: We analyze single-agent, adaptive, and multi-agent systems, emphasizing reasoning-centered workflow orchestration across in-context and post-training dimensions.


\item \textbf{Applications and evaluation}: We review real-world applications and benchmarks to illustrate the instantiation and evaluation of agentic reasoning mechanisms.
\item \textbf{Future agenda}: We identify emerging challenges in robustness, trustworthiness, and efficiency, outlining directions for the next generation of adaptive and collaborative agents.
\end{itemize}
\end{tcolorbox}

